% Evidence: Analyses003/004/025/028/043/051/054.
% Exact selections and provenance: revisions/2026-11-05-submission-synthesis/evidence_selection.json.
\section{Implications for Diagnosis and Intervention}
\label{sec:implications}

The preceding experiments locate cancellation and relate it to the observed
loss response. We now examine how these distinctions inform intervention
assessment. Language-model runs and matched changes to learning rate,
coordinate scaling and clipping provide complementary tests: an intervention
can alter gradient geometry, the optimizer's transformation, or both, while
its benefit depends on the objective and observation window.
Table~\ref{tab:submission-interventions} summarizes selected comparisons from
these studies. They illustrate different diagnostic questions; their effect
sizes are evaluated within each study and are not pooled.

\subsection{Language-Model Training and Optimizer Transformations}

We train randomly initialized Pythia configurations~\cite{biderman2023pythia}
with 14M, 31M and 70M parameters on MiniPile~\cite{kaddour2023minipile}.
The full-pass comparison uses three paired initialization/data-order seeds
per size and 712 AdamW updates. Each full update accumulates 1,024 sequences
of length 2,048, with microbatches of four; the final update includes the
remaining 311 sequences. Training uses token-weighted next-token
cross-entropy, AdamW coefficients $(0.9,0.95)$, $\epsilon=10^{-8}$,
global gradient clipping at one before the moment updates, and decay
0.1 excluding biases and LayerNorm parameters. The rate warms up linearly
over the first 1\% of the 712-update schedule to $10^{-3}$, then follows
cosine decay with a $10^{-4}$ floor. Parameters and moments
use FP32, training uses dynamically scaled FP16, and geometric reductions
use FP64. These raw-gradient measurements precede the transformations
that determine the applied update; their cumulative sum is not the
actual AdamW parameter displacement.

Across all three seed blocks, larger models attain lower validation loss
while exhibiting more frequent negative adjacent raw-gradient cosine.
Nevertheless, total raw-gradient interaction is positive in all nine runs.
Replacing all cross-step terms by the lag-one contribution underestimates
the squared cumulative gradient norm by 97.90--98.57\%. The distinction
between local cancellation and cumulative interaction thus also matters
in stochastic language-model training. Architecture and initialization
scales vary with model size, so this comparison does not identify a
dimensional mechanism for the loss ordering.

For a matched schedule comparison, the 14M and 70M models each branch
from their seed-1235 checkpoint after update 356. We compare the original
cosine decay, holding the current rate $5.1791\times10^{-4}$, and advancing
the decay clock twice as fast, all through update 712. Branches share
weights, moments, data position and future batches. Validation covers
all 500 documents and 693,329 supervised targets, including the short tail.
Holding the rate produces the most negative median cosine and the lowest
validation loss at both sizes; accelerated decay shifts cosine upward
but yields the highest loss. The schedule affects local geometry without
making positive alignment a performance target. Because the rate also
scales AdamW decay, this comparison concerns the complete schedule change.

An additional component intervention tests why the optimizer stream matters.
Across the same three sizes and three seed blocks, nine native trajectories
follow the first 350 updates of the original 712-update learning-rate clock
and supply matched forks after updates 175 and 225. Each of the
18 branches removes coordinate scaling for 50 updates: its own bias-corrected
first moment is normalized to the native parent's adaptive-step length
at each corresponding update. The branch continues updating its moments
and applies the same decay rule to its own weights. A fixed 128-sequence
probe, with 262,016 supervised targets, supplies a common token-weighted
objective. It is disjoint from the union of the three 350-update
training prefixes in this intervention study.
All 18 branches have higher endpoint probe loss than their native controls.
Table~\ref{tab:submission-interventions} reports the later fork's paired
geometric-mean penalties across seeds. Coordinate scaling changes the
finite-step response even when the adaptive-component length is matched.
This intervention isolates a specific direction substitution at the native
scale; it leaves the performance of independently tuned momentum methods
and mediation of the model-size gap open.

\begin{table*}[tp]
 \caption{Selected matched interventions. Each row uses its own control,
 objective and time window. In the first two rows, entries follow original
 decay / held rate / accelerated decay; cosine medians use affected pairs
 358--712 and validation loss is measured at update 712. Scaling-ablation
 penalties are $100[\exp(\operatorname{mean}_{\mathrm{seed}}
 \log(\mathcal L_{\mathrm{ablated}}/\mathcal L_{\mathrm{native}}))-1]$
 on the fixed probe. Other percentages compare the indicated training-loss
 statistic with its matched unchanged control; positive means higher loss.
 These selected development comparisons are not independent confirmations
 of a common intervention effect.}
 \label{tab:submission-interventions}
 \centering
 \footnotesize
 \setlength{\tabcolsep}{3pt}
 \renewcommand{\arraystretch}{1.12}
 % Generated values verified against revisions/2026-11-05-submission-synthesis/evidence_selection.json.
\begin{tabular}{>{\raggedright\arraybackslash}p{0.17\textwidth}>{\raggedright\arraybackslash}p{0.24\textwidth}>{\raggedright\arraybackslash}p{0.25\textwidth}>{\raggedright\arraybackslash}p{0.27\textwidth}}
\toprule
Case and replication & Matched comparison & Diagnostic or treatment response & Loss outcome \\
\midrule
Pythia-14M; one seed & Three schedules, updates 357--712
 & Median raw $\rho$: $0.194/-0.358/0.331$
 & Full validation: $5.1820/5.0674/5.2176$ \\
Pythia-70M; one seed & Same schedule comparison
 & Median raw $\rho$: $-0.109/-0.311/-0.006$
 & Full validation: $4.0559/3.9846/4.1042$ \\
Pythia-14M/31M/70M; three seeds
 & Remove coordinate scaling; updates 226--275
 & Native-parent adaptive-step norm matched at each update
 & Fixed-probe penalties: $26.84/57.90/106.49\%$ \\
\midrule
Transformer; one seed & Half LR; 64 updates from state 1,024
 & Mean raw $\rho$: $-0.299\to+0.589$
 & Mean $+3.23\%$; final $+4.86\%$ \\
MLP; one seed & Raw gradient cap 1; 288 updates from state 512
 & Clipping active on $125/288$ updates
 & Mean $+48.01\%$; maximum $3.82\times$; final $-19.26\%$ \\
ResNet-32; one seed & Raw gradient cap 1; 320 updates from state 256
 & Clipping active on $318/320$ updates
 & Mean $-28.32\%$; maximum $-52.15\%$; last-32 mean $+78.03\%$ \\
\bottomrule
\end{tabular}

\end{table*}

\subsection{The Response Window Changes the Intervention Assessment}

The full-objective interventions inherit the Adam states and objectives
of Section~\ref{sec:adam-diagnostics}: 5,000-image half-squared error for
the MLP, 50,000-image cross-entropy for ResNet, and 2,088,600-target
negative log-likelihood for the Transformer. All use seed zero,
$(\beta_1,\beta_2)=(0.9,0.999)$ and $\epsilon=10^{-8}$.
Treatment and control share a warm checkpoint and numerical environment,
with verified within-environment replay. We use fresh controls when
cross-device historical replay does not qualify. Learning rates remain
0.001 for the MLP and 0.003 for ResNet/Transformer unless explicitly changed;
the selected comparisons use no weight decay.

From Transformer state 1,024, halving the rate for 64 updates changes mean
raw cosine from $-0.299$ to $+0.589$, while mean and final training losses
become 3.23\% and 4.86\% higher than unchanged Adam. More positive
raw-gradient alignment therefore accompanies slower progress in this window.

Clipping supplies a different distinction: a cap on the gradient entering
Adam acts before the moments and coordinate scaling. We apply global
norm clipping at one over MLP states 512--800 and ResNet states 256--576,
with raw-gradient diagnostics measured before clipping. The cap activates
on 125/288 and 318/320 updates, respectively. The recorded parameter motion
shows why input clipping and step control differ: the MLP's mean actual
step norm increases from 0.04179 to 0.05707 and its maximum from 0.09366
to 0.18368, whereas ResNet's maximum decreases from 0.600 to 0.158.
The cap bounds the input to Adam; moment history and coordinate scaling
determine how that input contributes to parameter displacement.
In the MLP, the first peak
shrinks, but a later excursion makes the maximum loss 3.82 times the
control and raises mean loss by 48.01\%. Its final loss is nevertheless
19.26\% lower. In ResNet, clipping lowers window-average loss by 28.32\%
and the maximum by 52.15\%, but raises the last-32-state mean by 78.03\%.
Endpoint, average and excursion severity can thus support different
assessments of the same intervention. These examples concern selected
warm states; they establish neither a generally preferred clipping
threshold nor a uniform benefit from reducing step size.

\subsection{Using the Diagnostics and Interpreting Their Scope}

A constructive use begins by identifying the vector stream and the
question being asked. Raw gradients describe the changing local objective;
actual increments describe parameter motion. Adjacent cosine and norm
ratios locate reversal and amplification. The factors $C_T$ and $D_T$
separate cancellation within adjacent pairs from turning between their
residual displacements; the sign of $Q_T-1$ identifies the sign of the
all-pair interaction. Under adaptive optimization, alignment between
the current gradient and the actual increment additionally tests the
first-order loss contribution. Loss measured on a common objective then determines
whether the observed motion accompanies progress. An intervention can
be assessed against matched states with this same sequence of measurements,
including its activation, replacement excursions and later loss response.
The resulting diagnosis specifies what changed and which outcome it affected,
without requiring a forecasting rule.

The evidence has a defined scope. GD and full-batch Adam sweeps use one
initialization per model; the SGD and Pythia comparisons supply three
initialization blocks under their respective recipes. Selected intervention
windows and repeated events do not add independent seeds. Parameter-space
geometry depends on parameterization and does not directly measure change
in represented functions or stored knowledge. Observational model-size
comparisons, finite training horizons and same-corpus probes also constrain
generalization. The measurements provide a basis for formulating and testing
future intervention hypotheses; the present paper proposes no new training
method or validated predictive signal.
